\section{Mellin, Euler, convolution, Laplace, and transform pairs}\label{sec:transforms}

The source page explicitly omits many convergence restrictions.  We retain the
formulas but replace ``whenever the integral exists'' by checkable endpoint,
strip, branch, and Fubini certificates.

\subsection{A finite Mellin-pair theorem}

Write the plus-sign kernel as
\begin{equation}\label{eq:M-kernel}
 M(s)=
 \frac{\prod_{j=1}^{m}\Gamma(b_j+s)
       \prod_{j=1}^{n}\Gamma(1-a_j-s)}
      {\prod_{j=m+1}^{q}\Gamma(1-b_j-s)
       \prod_{j=n+1}^{p}\Gamma(a_j+s)}.
\end{equation}
The expected pair is
\begin{equation}\label{eq:mellin-pair}
 G(\eta x)=\frac1{2\pi\ii}\int_{c-\ii\infty}^{c+\ii\infty}
 M(s)(\eta x)^{-s}\dd s,
 \qquad
 \int_0^\infty x^{u-1}G(\eta x)\dd x=\eta^{-u}M(u).
\end{equation}
A raw exchange of these two infinite integrals would insert a Dirac delta and
is not a valid elementary Fubini proof.  The constructive proof uses finite
Fourier inversion after $x=\e^y$.

\begin{definition}[Mellin certificate]\label{def:mellin-cert}
A Mellin-pair certificate at $c$ gives:
\begin{enumerate}
\item endpoint majorants for $f(y)=\e^{cy}G(\eta\e^y)$ and its required finite
      differences as $y\to\pm\infty$;
\item vertical decay and quadrature bounds for $M(c+\ii t)$;
\item a rational modulus of continuity for $f$ on every requested finite
      interval;
\item an error schedule for the Fej\'er kernels
\[
 K_T(v)=\frac1{2\pi T}\left(\frac{\sin(Tv/2)}{\sin(v/2)}\right)^2
\]
(or the continuous triangular Fourier analogue).
\end{enumerate}
\end{definition}
All occurrences of trigonometric functions in the finite kernel have certified
power-series approximants; its positivity and normalization reduce to finite
identities and rational error bounds.

\begin{theorem}[Certified Mellin transform]\label{thm:mellin}
Assume a Mellin certificate and a rational strip
\begin{equation}\label{eq:mellin-strip}
 -\min_{1\le j\le m}\Re b_j+\epsilon
 \le\Re u\le
 1-\max_{1\le j\le n}\Re a_j-\epsilon
\end{equation}
with $\epsilon>0$ (with the absent side omitted when $m=0$ or $n=0$).  Then
\begin{equation}\label{eq:mellin-result}
 \int_0^\infty x^{u-1}G_{p,q}^{m,n}(\eta x)\dd x
 =\eta^{-u}
 \frac{\prod_{j=1}^{m}\Gamma(b_j+u)
       \prod_{j=1}^{n}\Gamma(1-a_j-u)}
      {\prod_{j=m+1}^{q}\Gamma(1-b_j-u)
       \prod_{j=n+1}^{p}\Gamma(a_j+u)}.
\end{equation}
The equality is a family of finite ball equalities with a primitive-recursive
modulus.
\end{theorem}
\begin{proof}
After $x=\e^y$, the Mellin transform on the line $\Re u=c$ is the Fourier
transform of $f(y)$.  Truncate $y$ and $t$ to rational intervals, use finite
quadrature, and replace the sharp cutoff by the triangular multiplier.  The
resulting double sum may be reordered exactly.  The finite Fourier kernel
approximates the identity by the supplied modulus of continuity, while endpoint
and frequency tails use the certificate.  Moving $u$ within the strip is
justified by multiplying $f(y)$ by $\e^{(\Re u-c)y}$; the margin $\epsilon$
preserves both tails.  The Gamma ratio is $M(u)$ by definition.  Every error is
assigned a dyadic budget before the finite computation.
\end{proof}

\subsection{Euler--Beta integrals}

Insert the finite contour representation of $G(zx)$ and use
\begin{equation}\label{eq:beta-inner}
 \int_0^1 x^{s-\alpha}(1-x)^{\alpha-\beta-1}\dd x
 =\frac{\Gamma(1-\alpha+s)\Gamma(\alpha-\beta)}
        {\Gamma(1-\beta+s)}.
\end{equation}
The new Gamma quotient is exactly one new top numerator parameter $\alpha$ and
one new bottom denominator parameter $\beta$.

\begin{theorem}[Two Euler-type integrals]\label{thm:euler-integrals}
Assume $\Re(\alpha-\beta)>0$, compatible branches, and an endpoint/Fubini
certificate.  Then
\begin{align}
\int_0^1x^{-\alpha}(1-x)^{\alpha-\beta-1}
G_{p,q}^{m,n}(zx)\dd x
&=\Gamma(\alpha-\beta)
G_{p+1,q+1}^{m,n+1}\!\left(
 \left.\begin{matrix}\alpha,\mathbf a\\\mathbf b,\beta\end{matrix}\right|z\right),
\label{eq:euler-zero-one}\\
\int_1^\infty x^{-\alpha}(x-1)^{\alpha-\beta-1}
G_{p,q}^{m,n}(zx)\dd x
&=\Gamma(\alpha-\beta)
G_{p+1,q+1}^{m+1,n}\!\left(
 \left.\begin{matrix}\mathbf a,\alpha\\\beta,\mathbf b\end{matrix}\right|z\right).
\label{eq:euler-one-infty}
\end{align}
\end{theorem}
\begin{proof}
For \eqref{eq:euler-zero-one}, perform the calculation
\eqref{eq:beta-inner} on $[\epsilon,1-\epsilon]$, where all sums and integrals
are finite.  The beta identity itself is proved by rational quadrature and the
Gamma/Beta contract.  The endpoint majorants send $\epsilon$ to the requested
precision.  For \eqref{eq:euler-one-infty}, substitute $x=1/t$; the inner
integral becomes a beta integral whose Gamma quotient is
$\Gamma(\beta-s)/\Gamma(\alpha-s)$, inserting the bottom numerator parameter
$\beta$ and top denominator parameter $\alpha$.  The supplied Fubini budget
justifies the finite order exchange.
\end{proof}

\subsection{Convolution of two \texorpdfstring{$G$}{G}-functions}

Let the second $G$ have parameters $\mathbf c$ (length $\sigma$),
$\mathbf d$ (length $\tau$), and indices $\mu,\nu$.  The source gives two
equivalent concatenations.

\begin{theorem}[Meijer convolution]\label{thm:convolution}
Under a common Mellin strip, positive scales $\eta,\omega$, and a convolution
certificate,
\begin{align}
&\int_0^\infty
G_{p,q}^{m,n}(\eta x)G_{\sigma,\tau}^{\mu,\nu}(\omega x)\dd x\notag\\
&=\frac1\eta
G_{q+\sigma,p+\tau}^{n+\mu,m+\nu}\!\left(
 \left.\begin{matrix}-b_1,\ldots,-b_m,\mathbf c,-b_{m+1},\ldots,-b_q\\
 -a_1,\ldots,-a_n,\mathbf d,-a_{n+1},\ldots,-a_p
 \end{matrix}\right|\frac\omega\eta\right)\label{eq:conv-one}\\
&=\frac1\omega
G_{p+\tau,q+\sigma}^{m+\nu,n+\mu}\!\left(
 \left.\begin{matrix}a_1,\ldots,a_n,-\mathbf d,a_{n+1},\ldots,a_p\\
 b_1,\ldots,b_m,-\mathbf c,b_{m+1},\ldots,b_q
 \end{matrix}\right|\frac\eta\omega\right).
\label{eq:conv-two}
\end{align}
\end{theorem}
\begin{proof}
Insert a finite Mellin--Barnes approximation for the first factor.  The inner
$x$-integral is the Mellin transform of the second factor at $1-s$, supplied by
\cref{thm:mellin}.  Its Gamma factors concatenate with those already present;
reading the four blocks gives \eqref{eq:conv-one}.  Insert the second factor
instead to obtain \eqref{eq:conv-two}; alternatively apply inversion and the
scale change $x\mapsto(\eta/\omega)x$.  The convolution certificate bounds the
four discarded regions in the finite $(x,t)$ rectangle, so the order exchange
is an iPRA finite-sum equality plus error inequalities.
\end{proof}

\subsection{Products with rational powers}

The introductory source claim allows arguments $cx^\gamma$ and $dx^\delta$.
The following finite compiler proves its mathematically precise form.

\begin{theorem}[Rational-power product compiler]\label{thm:rational-power-product}
Let $\gamma,\delta>0$ be rational, write
$\gamma/\delta=A/B$ in lowest terms with $A,B\in\N_{>0}$, and let $c,d>0$.
Under a common endpoint, branch, Mellin, and Fubini certificate, every integral
\begin{equation}\label{eq:rational-power-product}
 \int_0^\infty x^{\lambda-1}G_1(cx^\gamma)G_2(dx^\delta)\dd x
\end{equation}
is a computable constant times one ordinary Meijer $G$-function.  Its four
parameter blocks, argument, and constant are produced by a finite
primitive-recursive list transformation.  Taking $\lambda=1$ proves the
product-integral assertion on the source page.
\end{theorem}
\begin{proof}
Insert
$G_1(cx^\gamma)=(2\pi\ii)^{-1}\int\Phi_1(s)c^sx^{\gamma s}\dd s$.
The inner integral is the power-substituted Mellin transform
\[
 \int_0^\infty x^{\lambda+\gamma s-1}G_2(dx^\delta)\dd x
 =\frac1\delta d^{-(\lambda+\gamma s)/\delta}
 M_2\!\left(\frac{\lambda+\gamma s}{\delta}\right).
\]
This is first performed on a finite rectangle; the certificate controls all
discarded pieces.  Put $s=Bu$.  Gamma factors from $\Phi_1(Bu)$ have slope
$\pm B$, while those from $M_2(\lambda/\delta+Au)$ have slope $\pm A$.
Gauss multiplication replaces them by unit-slope factors, for example
\begin{align*}
 \Gamma(\alpha+Au)
 &= (2\pi)^{(1-A)/2}A^{\alpha+Au-1/2}
    \prod_{r=0}^{A-1}\Gamma\!\left(u+\frac{\alpha+r}{A}\right),\\
 \Gamma(\beta-Bu)
 &= (2\pi)^{(1-B)/2}B^{\beta-Bu-1/2}
    \prod_{r=0}^{B-1}\Gamma\!\left(\frac{\beta+r}{B}-u\right).
\end{align*}
The inverted formulas handle reciprocal Gamma factors.  Tagging each result by
slope sign and numerator/denominator position sorts the finite list into the
four standard blocks.  All $u$-dependent powers combine into one new argument;
the rest form the prefactor.  Hence the resulting integral is one ordinary
Meijer $G$-integral.  List expansion, sorting, affine parameter formation, and
exponent collection are bounded loops, and tail transport uses the same
finite Gamma-multiplication bounds as \cref{thm:multiplication}.
\end{proof}

\subsection{Laplace transforms}

The elementary inner integral is
\begin{equation}\label{eq:laplace-inner}
 \int_0^\infty \e^{-\omega x}x^{s-\alpha}\dd x
 =\omega^{\alpha-s-1}\Gamma(1-\alpha+s),
 \qquad \Re\omega>0.
\end{equation}

\begin{theorem}[Laplace and inverse Laplace]\label{thm:laplace}
Under endpoint and Fubini certificates and $\Re\omega>0$,
\begin{equation}\label{eq:laplace-main}
\int_0^\infty \e^{-\omega x}x^{-\alpha}G_{p,q}^{m,n}(\eta x)\dd x
=\omega^{\alpha-1}
G_{p+1,q}^{m,n+1}\!\left(
 \left.\begin{matrix}\alpha,\mathbf a\\\mathbf b\end{matrix}\right|
 \frac\eta\omega\right).
\end{equation}
If a Bromwich-line certificate places $\Re\omega=c$ to the right of all poles,
then the inverse formula is
\begin{equation}\label{eq:inverse-laplace}
 x^{-\alpha}G_{p,q+1}^{m,n}(\mathbf a;\mathbf b,\alpha\mid\eta x)
 =\frac1{2\pi\ii}\int_{c-\ii\infty}^{c+\ii\infty}
 \e^{\omega x}\omega^{\alpha-1}
 G_{p,q}^{m,n}(\mathbf a;\mathbf b\mid\eta/\omega)\dd\omega.
\end{equation}
\end{theorem}
\begin{proof}
Equation \eqref{eq:laplace-inner} inserts the top numerator parameter $\alpha$
and extracts $\omega^{\alpha-1}$; the remaining power changes the argument to
$\eta/\omega$.  Finite Fubini proves \eqref{eq:laplace-main}.  The inverse is
finite Bromwich/Fourier inversion after $\omega=c+\ii t$, with exactly the same
Fej\'er-kernel method as \cref{thm:mellin}.  The line and tail bounds are part
of the certificate, rather than hidden in the phrase ``as usual''.
\end{proof}

\begin{theorem}[Quadratic-argument Laplace formula]\label{thm:quadratic-laplace}
For $\Re\omega>0$ and the same type of certificates,
\begin{equation}\label{eq:quadratic-laplace}
 \int_0^\infty \e^{-\omega x}G_{p,q}^{m,n}(\eta x^2)\dd x
 =\frac1{\sqrt\pi\,\omega}
 G_{p+2,q}^{m,n+2}\!\left(
 \left.\begin{matrix}0,\frac12,\mathbf a\\\mathbf b\end{matrix}\right|
 \frac{4\eta}{\omega^2}\right).
\end{equation}
\end{theorem}
\begin{proof}
Inside a finite contour use
$\int_0^\infty\e^{-\omega x}x^{2s}\dd x=
\omega^{-2s-1}\Gamma(2s+1)$ and the duplication identity
\[
 \Gamma(2s+1)=\frac{2^{2s}}{\sqrt\pi}\Gamma(s+1)\Gamma(s+\tfrac12).
\]
The two new Gamma factors correspond to top parameters $0$ and $1/2$, and
$2^{2s}\eta^s\omega^{-2s}=(4\eta/\omega^2)^s$.
\end{proof}

\subsection{Transform kernels: what is proved}

A pair $(k,h)$ is a transform pair on a coded test class $\mathcal T$ when the
two finite-window operators
\[
 (K_Rf)(z)=\int_{1/R}^{R}k(z,y)f(y)\dd y,
 \qquad
 (H_Rg)(z)=\int_{1/R}^{R}h(z,y)g(y)\dd y
\]
have quadrature and tail certificates, and
$H_{R(k)}K_{R(k)}f$ is within $2^{-k}$ of $f$ on the requested rational compact
set.  This definition is first-order once $f,k,h,R$ are fixed PR names.

\begin{theorem}[Certified kernel principle]\label{thm:kernel-principle}
Suppose the Mellin multipliers of $k$ and $h$ multiply to $1$ on a rational
strip, and suppose the test function and kernels carry common Mellin, endpoint,
and Fubini certificates.  Then $(k,h)$ is a transform pair in the finite sense
above.
\end{theorem}
\begin{proof}
Take finite Mellin transforms.  Convolution becomes multiplication by
\cref{thm:convolution}; the two multipliers cancel algebraically.  Apply the
finite Mellin inversion theorem and add the four explicit truncation errors.
\end{proof}

The Narain kernels on the source page have reciprocal Gamma-ratio multipliers
once
$n-p=m-q>0$ and
$\sum a_j+\sum b_j=\sum c_j+\sum d_j$.  Therefore they instantiate
\cref{thm:kernel-principle} whenever the further endpoint conditions are
supplied.  If $p=q$, $m=n$, $a_j+b_j=0$, and $c_j+d_j=0$, parameter inversion
and reordering show $k=h$; the Hankel kernel is the stated special case.  The
Wimp kernels are treated identically after expanding the auxiliary $A$-function
into its $G$-kernel.  These are status \statusQ/\statusR results because the
source does not specify a maximal test class.

The generalized Laplace pair follows from \cref{thm:kernel-principle} with its
stated balance
\[
 (q+1-p)\frac{\rho}{2\gamma}=\sum_ja_j-\sum_jb_j,
 \qquad \gamma>0,\quad p\le q.
\]
At $\gamma=1/2$, $\rho=0$, and $p=q=0$, the kernel reduces by
\cref{thm:exp-case} to the ordinary Laplace kernel.  The two Meijer transforms
with $K_\nu/I_\nu$ and with Whittaker $W/M$ are direct specializations.  For the
second, set $\gamma=1/2$, $\rho=-m-k$, $p=q=1$, $a_1=m-k$, and $b_1=2m$.
The forward-kernel identification is the residue identity
\begin{equation}\label{eq:g-u-whittaker}
 G_{1,2}^{2,0}\!\left(
  \left.\begin{matrix}a\\b_1,b_2\end{matrix}\right|x\right)
 =x^{b_2}\e^{-x}U(a-b_1,1+b_2-b_1,x),
\end{equation}
with a common branch.  Taking $(a,b_1,b_2)=(m-k,0,2m)$ and using
\[
 W_{k+1/2,m}(x)=x^{m+1/2}\e^{-x/2}U(m-k,1+2m,x)
\]
gives exactly the source's $x^{-k-1/2}\e^{-x/2}W_{k+1/2,m}(x)$ kernel after
multiplying by the generalized-Laplace power $x^{-m-k}$.  For the inverse
kernel, \cref{thm:hyp-to-g} gives
\[
 G_{1,2}^{1,1}\!\left(
 \left.\begin{matrix}k-m\\0,-2m\end{matrix}\right|-x\right)
 =\frac{\Gamma(1+m-k)}{\Gamma(1+2m)}
 {}_1F_1(1+m-k;1+2m;x).
\]
Combining this with
$M_{k-1/2,m}(x)=x^{m+1/2}\e^{-x/2}
 {}_1F_1(1+m-k;1+2m;x)$ produces the displayed inverse kernel and its Gamma
normalization.  Formula \eqref{eq:g-u-whittaker} follows from the same two
residue families as the upper incomplete Gamma case, with a factorial tail.
Thus both Meijer pairs have status \statusQ under the inherited certificates.
An unrestricted theorem for every ``suitable'' function is intentionally not
asserted.
